% Vietnamese translation for OI Public Library
% Source-PDF: IOI中国国家候选队论文/2004/楼天城.pdf
\documentclass[11pt]{article}
\usepackage{oipl-vn}

\newfontfamily\cjkfont{WenQuanYi Zen Hei}
\newcommand{\zh}[1]{{\cjkfont #1}}
\newcommand{\Oh}{\mathcal{O}}
\newcommand{\blank}{\texttt{2}}

\title{Bàn về ứng dụng tìm kiếm bộ phận kết hợp thuật toán hiệu quả trong các bài toán tìm kiếm}
\author{Lou Tiancheng\\Trường Trung học số 14 Hàng Châu, Chiết Giang}
\date{2004}

\begin{document}
\maketitle

\origtitle{A brief discussion of partial search + efficient algorithms in search problems}
\sourcepath{IOI China candidate team papers / 2004 / Lou Tiancheng}

\begin{abstract}
Bài viết bắt đầu từ việc tìm kiếm trong các bài toán ghép cặp có ràng buộc
vị trí. Qua phân tích bài \textit{Milk Bottle Data}, bài viết đưa ra một kiểu
tìm kiếm không thông thường của tìm kiếm theo chiều sâu: \term{tìm kiếm bộ
phận + thuật toán hiệu quả}. Sau đó, thông qua ứng dụng của tìm kiếm bộ phận
trong hai bài \textit{Triangle Construction} và \textit{Phá trận liên hoàn}
(\zh{智破连环阵}), bài viết thảo luận các cách tối ưu chủ yếu dùng chung
cho phương pháp này, đồng thời phân tích bản chất của nó để chỉ ra vì sao nó
hiệu quả, cũng như các điều kiện và hạn chế cần thỏa khi áp dụng.
\end{abstract}

\noindent\textbf{Từ khóa:} tìm kiếm bộ phận; thuật toán hiệu quả.

\section{Mở đầu}

Trong nhiều bài toán, nếu có thể xây dựng mô hình toán học thì ta nên cố gắng
dùng phương pháp giải tích để xử lý, vì mô hình đơn giản phản ánh rõ hơn quan
hệ giữa các đối tượng.

Tuy nhiên, không phải bài toán nào cũng có thể xây dựng một mô hình toán học
đơn giản. Khi đó ta buộc phải dùng phương pháp tìm kiếm, tức liệt kê mọi khả
năng để tìm nghiệm khả thi hoặc nghiệm tối ưu.

Vì tìm kiếm dựa trên liệt kê, nên nó thường gắn với hiệu suất thấp. Có lúc
khối lượng tính toán của tìm kiếm quá lớn, khiến việc xử lý trở nên rất nặng.

Do đó ta cần dùng nhiều kỹ thuật để nâng cao hiệu suất. Cắt tỉa theo tính khả
thi, cắt tỉa theo tính tối ưu và điều chỉnh thứ tự tìm kiếm đều rất hữu ích;
nhờ chúng, hiệu suất tìm kiếm có thể tăng lên đáng kể.

Nhưng với một số bài, các tối ưu thông thường ấy khó có đất dụng võ. Khi đó ta
phải dùng các phương pháp tìm kiếm không thông thường. Trong bài này ta thảo
luận một trong số đó: \term{tìm kiếm bộ phận + thuật toán hiệu quả}.

\section{Bài dẫn}

Có $N$ vật phẩm và $N$ vị trí. Với mỗi vật phẩm, ta biết tập các vị trí mà nó
có thể đặt vào; cần tìm một song ánh giữa vật phẩm và vị trí. Ngoài ra còn có
các ràng buộc giữa vật phẩm và vị trí, chẳng hạn: nếu vật phẩm $1$ đặt ở vị
trí $3$ thì vật phẩm $2$ không được đặt ở vị trí $1$. Yêu cầu có thể là tìm
một nghiệm khả thi, hoặc gán trọng số cho mỗi cách tương ứng rồi tìm nghiệm
tối ưu thỏa điều kiện.

Do quan hệ ràng buộc giữa các đối tượng rất phức tạp, rất khó xây dựng một
quan hệ đồ thị hai phía đơn giản, hoặc dùng luồng mạng để giải.

Trước một loạt bài toán kiểu này, thông thường ta chỉ còn cách tìm kiếm. Vấn
đề là tìm kiếm thế nào và tối ưu ra sao.

\subsection{Phân tích đơn giản}

Nếu liệt kê vị trí của từng vật phẩm rồi kiểm tra, độ phức tạp thời gian là
$\Oh(N!)$. Thoạt nhìn dường như cũng chỉ có thể làm như vậy.

\subsection{Phân tích thêm}

Xét một ví dụ với $N=6$. Có bốn ràng buộc khác, mỗi ràng buộc dạng
$(a,b,c,d)$ nghĩa là nếu $a$ đặt ở $b$ thì $c$ không được đặt ở $d$:
\begin{verbatim}
1 3 5 6
2 2 5 3
3 1 4 1
3 2 6 2
\end{verbatim}

Sau đó ta nhận ra: một khi đã xác định vị trí của vật phẩm $3$ và $5$, giữa vị
trí của bốn vật phẩm còn lại không còn quan hệ ràng buộc nào nữa. Khi đó vị
trí của bốn vật phẩm còn lại có thể giải bằng ghép cặp.

Từ đây ta thấy một mô hình tìm kiếm mới: \term{tìm kiếm bộ phận}.

\begin{quote}
\textbf{Tìm kiếm bộ phận}: tìm kiếm trên một phần biến để quan hệ giữa các
biến còn lại được đơn giản hóa; sau đó dùng một số thuật toán hiệu quả, thường
là ghép cặp, giải hệ phương trình, tham lam, quy hoạch động, v.v., để hoàn
thành phần còn lại của bài toán.
\end{quote}

Với bài dẫn trên, ta tìm kiếm trước vị trí của một số lượng vật phẩm nhất định
(không phải toàn bộ), làm cho quan hệ giữa các vật phẩm trong phần còn lại trở
thành quan hệ đồ thị hai phía, rồi dùng ghép cặp hai phía để xử lý các vật
phẩm còn lại.

\subsection{Bản chất}

Trong ví dụ trên, nếu sau khi biết vị trí của $3$ và $5$ mà ta không dùng ghép
cặp, thì thực chất ta đang dùng tìm kiếm để tìm một ghép cặp; hiệu suất tất
nhiên sẽ không cao.

Tìm kiếm bộ phận tạo điều kiện cho các thuật toán hiệu quả khác, chẳng hạn
tạo ra quan hệ hai phía trong ví dụ trên; còn thuật toán hiệu quả thay thế tìm
kiếm để hoàn thành nhanh phần việc còn lại.

Phương pháp tìm kiếm bộ phận phát huy đầy đủ ưu thế của cả tìm kiếm lẫn các
thuật toán hiệu quả khác. Ưu thế của tìm kiếm là phạm vi áp dụng rộng, có thể
vượt qua tình huống phức tạp; ưu thế của các thuật toán khác là hiệu suất cao.
Hai phía thúc đẩy nhau, đồng thời bù đắp nhược điểm của nhau. Đây chính là
chìa khóa thành công của phương pháp.

Tìm kiếm bộ phận kết hợp các thuật toán hiệu quả khác đã được dùng trong rất
nhiều bài. Dưới đây ta dùng vài ví dụ để thảo luận cách áp dụng và các kỹ
thuật tối ưu của kiểu tìm kiếm này.

\section{Ví dụ: Milk Bottle Data}

\subsection{Mô tả bài toán}

\textit{Milk Bottle Data} xuất hiện tại ACM/ICPC Asia Regional Shanghai 1996.

Một hộp được chia thành lưới $N\times N$; mỗi ô có thể chứa một chai sữa hoặc
không chứa gì. Ông Smith ghi lại tình trạng sữa của từng hàng từ trái sang
phải, và cũng ghi lại tình trạng sữa của từng cột từ trên xuống dưới. Mỗi bản
ghi gồm $N$ chữ số: $0$ nghĩa là không có sữa, $1$ nghĩa là có sữa. Không may,
thứ tự của $2N$ bản ghi đã bị xáo trộn, và một vài chữ số cũng bị mờ.

Bây giờ ông Smith nhờ bạn khôi phục tình trạng sữa ban đầu trong hộp.

\paragraph{Dữ liệu vào.}
Dòng đầu là số nguyên $N$. Tiếp theo là $2N$ dòng, mỗi dòng có $N$ chữ số:
$0$ nghĩa là chắc chắn không có sữa, $1$ nghĩa là chắc chắn có sữa, $2$ nghĩa
là không xác định.

\paragraph{Ví dụ.}
\begin{verbatim}
input
5
01210
21120
21001
12110
12101
12101
00011
22222
11001
10010
output
9 8 6 2 7
4 1 0 1 1 0
10 1 0 0 1 0
1 0 1 1 1 0
3 0 1 0 0 1
5 1 1 1 0 1
\end{verbatim}

Ràng buộc dữ liệu: $1\le N\le 10$.

\subsection{Phân tích ban đầu}

Quan hệ ràng buộc giữa hàng và cột rất phức tạp, khó tìm được thuật toán đa
thức. Luồng mạng cũng không dễ áp dụng.

Ta có thể dùng tìm kiếm theo chiều sâu với $(2N)!$ phương án: lần lượt liệt kê
số hiệu bản ghi cho từng hàng và từng cột, rồi kiểm tra có sinh mâu thuẫn hay
không.

Gọi chữ số thứ $j$ của bản ghi thứ $i$ là $A[i,j]$. Nếu hàng $a$ chọn bản ghi
$x$, cột $b$ chọn bản ghi $y$, thì điều kiện mâu thuẫn là
\[
  (A[x,b]\ne 2)\land(A[y,a]\ne 2)\land(A[x,b]\ne A[y,a]).
\]

\subsection{Phân tích hiệu suất thứ nhất}

Vì $N\le 10$, giá trị $(20)!\cdot 10$ lớn tới
\[
  24329020081766400000.
\]
Thử vài tối ưu cơ bản:
\begin{itemize}
  \item Cắt tỉa khả thi: ngoài kiểm tra sớm, gần như không còn tối ưu nào khác.
  \item Cắt tỉa tối ưu: bài toán không phải bài tối ưu.
  \item Điều chỉnh thứ tự tìm kiếm: nếu mỗi lần tìm kiếm nhánh có ít khả năng
  nhất thì hiệu quả vẫn khá tốt. Nhưng trường hợp xấu nhất vẫn phải tính
  $(20)!$, khó gọi là một thuật toán khả thi tốt.
\end{itemize}

\subsection{Tìm kiếm bộ phận + ghép cặp}

Ta nhận thấy giữa các hàng với nhau và giữa các cột với nhau không có ràng
buộc; mọi ràng buộc đều nằm giữa hàng và cột. Vậy ta có thể chỉ tìm kiếm các
hàng không? Câu trả lời là có.

Nếu đã biết giá trị của từng hàng, thì giữa các cột không còn quan hệ ràng
buộc. Khi đó ta có thể dùng ghép cặp để tìm một nghiệm khả thi.

Cách dựng đồ thị như sau. Phía trái của đồ thị hai phía có $N$ đỉnh, biểu
diễn $N$ bản ghi độ dài $N$ chưa được chọn; phía phải có $N$ đỉnh, biểu diễn
$N$ cột. Nếu đưa bản ghi ứng với đỉnh trái thứ $i$ vào cột $j$ không gây mâu
thuẫn ở bất kỳ hàng nào của cột $j$, thì thêm một cạnh giữa đỉnh trái thứ $i$
và đỉnh phải thứ $j$.

Ghép cặp hai phía giải được trong $\Oh(N^3)$, hiệu quả hơn nhiều so với tìm
kiếm mù.

\subsection{Phân tích hiệu suất thứ hai}

Khi $N=10$, khối lượng tính toán trở thành
\[
  P(20,10)\cdot(20\cdot 10\cdot 10)=13408850145600000.
\]
So với tìm kiếm đơn giản, hiệu suất đã được cải thiện rất nhiều. Vì chỉ cần
tìm một nghiệm là chương trình có thể kết thúc, cách này có thể nhanh chóng
vượt qua toàn bộ dữ liệu kiểm thử.

Ta có thể so sánh hai thuật toán như sau. Nếu tìm kiếm thông thường là tìm
kiếm hàng trước rồi tìm kiếm cột, thì sau khi tìm kiếm xong hàng, nó đang dùng
tìm kiếm để tìm một ghép cặp. Độ phức tạp của tìm kiếm là $\Oh(N!)$, còn độ
phức tạp của ghép cặp là $\Oh(N^3)$. Hiệu suất khác nhau là hiển nhiên.

Thuật toán tìm kiếm này vẫn có thể dùng các tối ưu thông thường. Phương pháp
tìm kiếm đặc biệt và các tối ưu thông thường không mâu thuẫn với nhau. Thuật
toán này còn có nhiều tối ưu riêng; ta tiếp tục qua hai ví dụ.

\section{Ví dụ: Triangle Construction}

\subsection{Mô tả bài toán}

\textit{Triangle Construction} là bài trong ZJU Monthly Contest.

Một tam giác đều cạnh $N$ có thể chia thành $N^2$ tam giác đều nhỏ cạnh $1$.
Trong hình gốc, các tam giác nhỏ được ghép lại sao cho hai cạnh kề nhau phải
có cùng màu. Bài toán cho màu của ba cạnh của từng tam giác trong $N^2$ tam
giác nhỏ; hãy xác định có thể dùng chúng để tạo thành một tam giác đều cạnh
$N$ hay không.

Ở đây bài viết chỉ xét trường hợp $N=4$.

\subsection{Phân tích}

Tìm kiếm trực tiếp có khối lượng $16!=20922789888000$, quá lớn.

Thực ra đây cũng là một bài ghép cặp có ràng buộc vị trí với $16$ vật thể.
Khác với bài dẫn ở chỗ: ràng buộc giữa các vị trí là xác định, vì vậy ta phải
tận dụng đầy đủ đặc điểm này.

Nếu dùng thuật toán tìm kiếm bộ phận + ghép cặp, ta nên tìm kiếm thế nào?
Đương nhiên ta cần chọn cách có hiệu suất cao nhất.

Giả sử số tam giác được liệt kê là $T$, khối lượng tính toán là
\[
  P(16,T)\cdot 256\cdot(16-T).
\]
Các giá trị trong bài gốc được liệt kê như sau:

\begin{center}
\begin{tabular}{r|r@{\qquad}r|r}
$T$ & Khối lượng & $T$ & Khối lượng\\
\hline
0 & 4096 & 9 & 7439214182400\\
1 & 61440 & 10 & 44635285094400\\
2 & 860160 & 11 & 223176425472000\\
3 & 11182080 & 12 & 892705701888000\\
4 & 134184960 & 13 & 2678117105664000\\
5 & 1476034560 & 14 & 5356234211328000\\
6 & 14760345600 & 15 & 5356234211328000\\
7 & 132843110400 & 16 & ?\\
8 & 1062744883200 & & \\
\end{tabular}
\end{center}

Kết luận: khối lượng tính toán tăng dần theo $T$. Với dạng bài này, trong
trường hợp thông thường, số biến liệt kê càng ít càng tốt.

\paragraph{Tối ưu.}
Chọn các biến liệt kê tốt, sao cho số biến cần liệt kê ít nhất có thể.

Với bài này, ta có thể liệt kê sáu tam giác ở phần được tô bóng trong hình
gốc; mười tam giác còn lại đều có thể dựng thành quan hệ hai phía đơn giản.
Khối lượng tính toán là
\[
  P(16,6)\cdot 256\cdot 10=14760345600.
\]
Cộng thêm các tối ưu thông thường, hiệu suất được nâng lên rất nhiều.

\subsection{Nhận xét}

Khi chọn biến tìm kiếm, nguyên tắc chung là làm cho số biến tìm kiếm ít nhất
có thể, vì thuật toán hiệu quả có độ phức tạp đa thức, còn tìm kiếm là dạng
NP.

Cách chọn đôi khi có thể dựa trên phân tích, như hai bài vừa nêu; nhưng trong
một số bài, chương trình cần tự xác định biến nào phải tìm kiếm.

\section{Ví dụ: Phá trận liên hoàn (NOI 2003)}

Đây là một ví dụ kinh điển về tối ưu phương pháp tìm kiếm bộ phận.

\subsection{Mô tả bài toán}

Sau khi tiêu tốn hàng chục tỷ, nước B cuối cùng đã chế tạo ra vũ khí mới:
Zenith Protected Linked Hybrid Zone, gọi là \textit{trận liên hoàn}, và tuyên
bố đây là một vũ khí tự phát thông minh không thể đánh bại. Nhưng sau khi
trinh sát, nước A phát hiện trận liên hoàn thực ra gồm $M$ vũ khí độc lập.
Các vũ khí này được đánh số $1,2,\ldots,M$. Mỗi vũ khí có hai trạng thái:
trạng thái tự vệ vô địch và trạng thái tấn công. Ban đầu, vũ khí số $1$ ở
trạng thái tấn công, các vũ khí khác đều ở trạng thái tự vệ vô địch. Về sau,
mỗi khi vũ khí số $i$ ($1\le i<M$) bị tiêu diệt, sau $1$ giây vũ khí số
$i+1$ tự động chuyển từ trạng thái tự vệ vô địch sang trạng thái tấn công.
Khi vũ khí số $M$ bị tiêu diệt, trận liên hoàn đắt đỏ này bị phá hủy hoàn
toàn.

Để đánh bại nước B, Bộ trưởng quân sự nước A định dùng loại vũ khí rẻ nhất:
bom. Sau thời gian dài thăm dò chính xác, các nhà quân sự nước A nắm được tọa
độ phẳng của $M$ vũ khí trong trận liên hoàn; dựa vào đó họ chọn $n$ điểm và
đặt bom hẹn giờ đặc biệt tại các điểm này. Các quả bom được đánh số
$1,2,\ldots,n$. Bán kính tác dụng của mỗi quả bom đều là $k$, và nó nổ liên
tục trong $5$ phút. Trong $5$ phút đó, mỗi quả bom có thể tức thời tiêu diệt
bất kỳ số vũ khí B nào đang ở trạng thái tấn công và có khoảng cách đường
thẳng đến nó không vượt quá $k$. Tương tự trận liên hoàn, ban đầu bom số
$a_1$ nổ liên tục trong $5$ phút, sau đó bom số $a_2$ nổ liên tục trong
$5$ phút, tiếp theo là $a_3$, v.v., cho đến khi trận liên hoàn bị phá hủy.
Khi một quả bom nổ, các quả bom chưa nổ khác đều còn được che giấu dưới đất
và không bị bom của phe mình phá hủy.

Rõ ràng việc chọn $a_1,a_2,a_3,\ldots$ rất quan trọng. Một dãy tốt có thể phá
hủy trận liên hoàn chỉ với ít bom; một dãy xấu có thể dùng hết mọi quả bom mà
vẫn chưa phá hủy được. Hãy quyết định một dãy $a_1,a_2,a_3,\ldots$ sao cho
trận liên hoàn bị phá hủy trong thời gian bom số $a_x$ nổ. Giá trị $x$ cần
càng nhỏ càng tốt.

\paragraph{Dữ liệu vào.}
Tệp \texttt{zplhz.in} có dòng đầu gồm ba số nguyên $M,n,k$
($1\le M,n\le 100$, $1\le k\le 1000$), lần lượt là số vũ khí của nước B, số
bom của nước A, và phạm vi tấn công của bom. Tiếp theo là $M$ dòng, mỗi dòng
gồm một cặp số nguyên $x_i,y_i$ ($0\le x_i,y_i\le 10000$), là tọa độ của vũ
khí số $i$. Sau đó là $n$ dòng, mỗi dòng gồm một cặp số nguyên $u_i,v_i$
($0\le u_i,v_i\le 10000$), là tọa độ của bom số $i$. Dữ liệu vào bảo đảm hợp
lệ và có nghiệm.

Trong dữ liệu kiểm thử, các giá trị $x_i,y_i,u_i,v_i$ được sinh ngẫu nhiên.

\paragraph{Dữ liệu ra.}
Tệp \texttt{zplhz.out} có dòng đầu là một số nguyên $x$, tức số bom thực sự
được dùng. Dòng thứ hai gồm $x$ số nguyên, lần lượt là $a_1,a_2,\ldots,a_x$.

\subsection{Phân tích ban đầu}

Bom $I$ của nước A bắn trúng vũ khí $J$ của nước B khi và chỉ khi
\[
  (u[I]-x[J])^2+(v[I]-y[J])^2\le r^2.
\]
Khó tìm được một thuật toán đa thức để tìm nghiệm tối ưu. Với dạng bài này,
thông thường ta chỉ có chiến lược tìm kiếm.

Phân tích thêm, mỗi quả bom nhất định tiêu diệt một đoạn liên tiếp theo chỉ
số trong các vũ khí của nước B. Thời gian $5$ phút chỉ nói rằng một quả bom có
thể tiêu diệt tùy ý nhiều vũ khí B có chỉ số liên tiếp.

\subsection{Tìm kiếm bộ phận}

Để dùng tìm kiếm bộ phận trong bài này, ta cần một phép chuyển hóa: giả sử đã
chia các vũ khí B theo chỉ số thành $x$ đoạn $[S_i,T_i]$, trong đó
$S_1=1$, $T_i\ge S_i$ và $S_{i+1}=T_i+1$. Sau đó kiểm tra liệu có thể chọn
$x$ quả bom trong $N$ quả bom của nước A, mỗi quả tiêu diệt một đoạn tương
ứng hay không.

Thực chất ta chia tìm kiếm thành hai phần: trước tiên tìm kiếm cách chia vũ
khí B theo chỉ số thành $x$ đoạn; sau đó tìm kiếm xem có thể chọn $x$ quả bom
từ $N$ quả bom của nước A, mỗi quả tiêu diệt một đoạn, hay không. Phần thứ hai
có thể giải bằng ghép cặp.

Đặt $C[S][T][I]$ biểu diễn bom $I$ của nước A có thể tiêu diệt toàn bộ vũ khí
B từ $S,S+1,\ldots,T-1,T$ hay không. Khi đó:
\[
\begin{aligned}
  C[S][S][I] &=
    \bigl((u[I]-x[S])^2+(v[I]-y[S])^2\le R^2\bigr),\\
  C[S][T][I] &= C[S][T-1][I]\land C[T][T][I]\qquad(S<T).
\end{aligned}
\]
Tính $C$ mất $\Oh(N^3)$.

Dựng đồ thị: phía trái có $x$ đỉnh, biểu diễn $x$ đoạn của vũ khí B sau khi
chia; phía phải có $N$ đỉnh, biểu diễn $N$ quả bom của nước A. Có cạnh từ đỉnh
trái thứ $i$ đến đỉnh phải thứ $j$ khi $C[S_i][T_i][j]$ đúng.

Nhiệm vụ của tìm kiếm là chia vũ khí B theo chỉ số thành một số đoạn, rồi dùng
ghép cặp hai phía để kiểm tra.

\subsection{Phân tích hiệu suất thứ nhất}

Khung tìm kiếm cơ bản đã có. Tuy dữ liệu được sinh ngẫu nhiên, số cách chia
$M$ vũ khí B vẫn rất nhiều, đôi khi có thể cao tới $2^m$. Thời gian khó chịu
đựng; nếu dùng giới hạn thời gian cứng thì tính đúng đắn bị ảnh hưởng, hiệu
quả sẽ không tốt.

Chỉ có $4$ bộ dữ liệu có thể ra nghiệm trong giới hạn thời gian; với $6$ bộ
còn lại, nếu dùng giới hạn thời gian thì có thêm $1$ bộ có thể nhận được nghiệm
tối ưu.

\subsection{Các tối ưu đầu tiên}

Có thể phân tích tối ưu từ hai mặt: tính khả thi và tính tối ưu.

\paragraph{Tối ưu 1: cắt tỉa tối ưu.}
Nếu bom của nước A được phép dùng lặp lại, đặt
\[
  \mathit{Dist}[i]=\text{số bom ít nhất cần dùng để tiêu diệt vũ khí B từ }i
  \text{ đến }m.
\]
Giá trị này có thể tính bằng quy hoạch động:
\[
  \mathit{Dist}[m+1]=0,
\]
\[
  \mathit{Dist}[i]=
  \min\{\mathit{Dist}[j]+1\mid
       \mathit{Can}[i][j-1][k],\ 1\le k\le n,\ i<j\le n+1\}.
\]
Tính $\mathit{Dist}$ mất $\Oh(N^3)$. Từ đó có điều kiện cắt tỉa tối ưu:
\[
\text{nếu số bom đã dùng}+\mathit{Dist}[\text{vũ khí B tiếp theo}]
\ge \text{nghiệm tốt nhất hiện có, thì cắt tỉa.}
\]

\paragraph{Tối ưu 2: cắt tỉa khả thi.}
Với phương pháp tìm kiếm này, thường có hai tối ưu khả thi rất hiệu quả:
\begin{enumerate}
  \item Kiểm tra sớm xem ghép cặp có thể thành công hay không, tránh tìm kiếm
  thừa.
  \item Mỗi lần ghép cặp có thể mở rộng từ ghép cặp trước đó, không cần bắt
  đầu lại từ đầu.
\end{enumerate}

Nếu cách chia hiện tại đã không thể ghép cặp thành công, không cần tiếp tục
tìm kiếm. Chỉ cần mỗi khi tìm kiếm thêm một đoạn mới thì lập tức dùng ghép cặp
để kiểm tra. Mỗi lần tìm ghép cặp chỉ cần mở rộng từ nền tảng cũ.

Sau hai tối ưu trên, hiệu suất chương trình tăng đáng kể. Trong $10$ bộ kiểm
thử, có $8$ bộ ra nghiệm trong giới hạn thời gian; với $2$ bộ còn lại, nếu
dùng giới hạn thời gian thì có $1$ bộ có thể nhận được nghiệm tối ưu.

\subsection{Tối ưu thêm}

Tối ưu 2 tuy loại bỏ nhiều cách chia không cần thiết, nhưng lãng phí không ít
thời gian khi kiểm tra. Vì vậy, khi liệt kê độ dài đoạn chia tiếp theo, ta có
thể dựa trên cách chia và tình trạng ghép cặp trước đó, tức các cạnh đã được
ghép, để dùng tìm kiếm theo chiều rộng $\Oh(n^2)$ tính độ dài lớn nhất
$\mathit{maxL}$ của đoạn kế tiếp. Rõ ràng mọi độ dài của đoạn kế tiếp trong
$[1,\mathit{maxL}]$ đều chắc chắn tìm được một ghép cặp khả thi. Như vậy vừa
tiết kiệm thời gian kiểm tra, vừa có thể liệt kê độ dài đoạn từ dài đến ngắn,
giúp chương trình nhanh chóng tiến gần nghiệm tối ưu và đồng thời tăng hiệu
quả của điều kiện cắt tỉa 1.

Trước hết cần tính $\mathit{MaxT}$:
\[
  \mathit{MaxT}[i][S]=\text{chỉ số lớn nhất mà bom }i\text{ có thể tiêu diệt
  nếu bắt đầu từ }S.
\]
Nếu bom $i$ không tiêu diệt được $S$, thì $\mathit{MaxT}[i][S]=S-1$.

Ta có thể dùng quy hoạch động để tính $\mathit{MaxT}[i][S]$:
\[
  \mathit{MaxT}[i][S]=
  \begin{cases}
    S-1, & \text{nếu bom }i\text{ không tiêu diệt được }S,\\
    \mathit{MaxT}[i][S+1], & \text{nếu bom }i\text{ tiêu diệt được }S.
  \end{cases}
\]
và $\mathit{MaxT}[i][m+1]=m$. Thời gian tính $\mathit{MaxT}$ là $\Oh(N^2)$.

Cách cài đặt cụ thể xét $n$ đỉnh phía phải của đồ thị hai phía, tức $n$ quả
bom. Nếu đỉnh $i$ chưa được ghép, ta xem $i$ là có thể sử dụng. Với một đỉnh
$i$ đã được ghép, nếu tồn tại một đường xen kẽ đi từ một đỉnh chưa ghép nào
đó đến $i$, và $i$ là đỉnh ngoài trên đường ấy, thì $i$ cũng được xem là có
thể sử dụng.

Ta cần tính các đỉnh đã ghép mà những đường xen kẽ xuất phát từ đỉnh chưa ghép
có thể đạt tới. Chỉ cần từ mỗi đỉnh chưa ghép chạy tìm kiếm theo chiều rộng,
mất $\Oh(N^2)$. Cần chú ý tránh lặp: nếu một đỉnh đã ghép đã được xác định là
có thể sử dụng, thì không cần mở rộng nó lần nữa, vì khi xác định nó là đỉnh
có thể sử dụng ta đã mở rộng từ nó; mở rộng lại chỉ tạo ra lặp thừa.

Do đó
\[
  \mathit{MaxL}=\max\{\mathit{MaxT}[i][S]\mid i\text{ có thể sử dụng}\}.
\]

\subsection{Phân tích hiệu suất thứ ba}

Sau các tối ưu trên, toàn bộ dữ liệu đều ra nghiệm ngay lập tức, và mọi kết
quả đều là nghiệm tối ưu. Ngay cả với dữ liệu ngẫu nhiên $n=200$, chương trình
cũng có thể ra nghiệm tức thời; điều này cho thấy hiệu suất đã được cải thiện
rất nhiều.

\subsection{Tinh chỉnh thêm}

Còn hai tối ưu nữa, đôi khi hiệu quả không tốt nhưng vẫn đáng nhắc tới.

\paragraph{Tối ưu 3: phân nhánh cận.}
Cách này có thể tăng hiệu quả của điều kiện cắt tỉa 1, nhưng khi nghiệm tối ưu
khác xa $\mathit{Dist}[1]$, nó sẽ lãng phí một lượng thời gian nhất định.

\paragraph{Tối ưu 4: cải thiện giá trị \(\mathit{Dist}[i]\).}
Trong tối ưu 1, giá trị $\mathit{Dist}[i]$ đôi khi không tối ưu. Qua thử
nghiệm, nếu $\mathit{Dist}[i]$ lệch nghiệm tối ưu $1$ đơn vị, đặc biệt khi
$i$ nhỏ, tốc độ chương trình bị ảnh hưởng rõ rệt. Vì vậy có thể dùng cùng
kiểu tìm kiếm để tính $\mathit{Dist}[i]$; đáp án của bài chính là
$\mathit{Dist}[1]$. Cách này tăng hiệu quả của điều kiện cắt tỉa 1, nhưng đồng
thời việc tìm kiếm cho mỗi $i$ cũng tốn thời gian.

\subsection{So sánh kết quả chương trình}

\begin{center}
\begin{tabular}{l|cccccccccc}
 & 1 & 2 & 3 & 4 & 5 & 6 & 7 & 8 & 9 & 10\\
\hline
Tìm kiếm đơn giản
 & 0.00 & 0.00 & 0.50 & T.O. & 0.65 & T.O. & T.O. & T.O. & T.O. & T.O.\\
Tìm kiếm tối ưu
 & 0.01 & 0.01 & 0.10 & T.O. & 0.50 & 0.80 & T.O. & 0.55 & 0.30 & 0.80\\
Tối ưu thêm
 & 0.01 & 0.01 & 0.02 & 0.03 & 0.00 & 0.02 & 0.02 & 0.01 & 0.02 & 0.02\\
\end{tabular}
\end{center}

Với phương pháp tìm kiếm này, nói chung có hai tối ưu khả thi rất hiệu quả:
\begin{enumerate}
  \item Kiểm tra sớm xem có thể thành công hay không, tránh tìm kiếm thừa.
  \item Mỗi lần kiểm tra cố gắng tận dụng nhiều nhất kết quả kiểm tra trước đó.
\end{enumerate}

\section{Tổng kết}

Các ví dụ trong bài đều có thể áp dụng phương pháp tìm kiếm bộ phận để giải
hiệu quả, và chúng có những điểm chung rõ rệt về mặt tư tưởng. Quá trình suy
nghĩ tổng quát có thể mô tả bằng lời như sau:

\begin{center}
\begin{tabular}{c}
Khó nghĩ ra thuật toán đa thức; tìm kiếm thông thường đơn giản không giải
quyết được bài toán\\
$\downarrow$\\
Muốn giảm lượng tìm kiếm\\
$\downarrow$\\
Chọn các biến cần tìm kiếm\\
$\downarrow$\\
Xét các biến còn lại có thể giải bằng thuật toán hiệu quả hay không\\
\quad Nếu không, quay lại chọn biến tìm kiếm\\
$\downarrow$\\
Xét liệu thuật toán có nâng cao hiệu suất hay không\\
\quad Nếu không, quay lại chọn biến tìm kiếm\\
$\downarrow$\\
Thông qua tối ưu thu được một phương pháp tìm kiếm hiệu quả
\end{tabular}
\end{center}

Các tối ưu thông thường gồm:
\begin{enumerate}
  \item Chọn tốt các biến cần liệt kê, làm cho số biến liệt kê ít nhất có thể.
  \item Tận dụng đầy đủ nhất kết quả trước đó, kiểm tra sớm và tránh tìm kiếm
  thừa.
\end{enumerate}

Tìm kiếm bộ phận cũng có thể kết hợp với giải hệ phương trình, tham lam, quy
hoạch động và các thuật toán hiệu quả khác.

Tìm kiếm bộ phận + thuật toán hiệu quả thể hiện sự kết hợp hữu cơ giữa tìm
kiếm và các phương pháp khác, phát huy đầy đủ ưu điểm của cả hai và bù đắp
nhược điểm của nhau. Đây là nguyên nhân chính tạo nên hiệu quả của nó. Vì vậy,
trong các bài toán tìm kiếm, việc linh hoạt áp dụng tìm kiếm bộ phận thường
có thể tạo ra hiệu quả bất ngờ.

Cần chú ý rằng dùng tìm kiếm bộ phận để giải bài toán tìm kiếm là một phương
pháp tìm kiếm không thông thường. Tuy trong các ví dụ của bài, tìm kiếm bộ
phận có nhiều ưu điểm vượt trội, nhưng không thể cho rằng phương pháp thông
thường nhất định kém hơn phương pháp không thông thường. Đa số bài toán tìm
kiếm vẫn phù hợp với các phương pháp tìm kiếm thông thường. Chỉ khi nắm vững
đặc điểm của tìm kiếm bộ phận và kết hợp nó nhuần nhuyễn với tìm kiếm thông
thường, ta mới thật sự thu được một thuật toán tìm kiếm hiệu quả.

\section*{Tài liệu tham khảo}

\begin{itemize}
  \item ACM/ICPC Asia Regional Shanghai 1996, \textit{Milk Bottle Data}.
  \item ZJU Monthly Contest, \textit{Triangle Construction}.
  \item NOI 2003, \textit{Phá trận liên hoàn} (\zh{智破连环阵}).
\end{itemize}

\section*{Lời cảm ơn}

Cảm ơn Zhejiang University Online Judge đã cung cấp mô tả bài
\textit{Milk Bottle Data} và \textit{Triangle Construction}.

\appendix
\section{Chương trình tham khảo cho Phá trận liên hoàn}

Phần phụ lục gốc đưa ra một chương trình C/C++ tham khảo. Các chú thích trong
mã dưới đây được dịch sang tiếng Việt; cấu trúc chương trình và tên biến được
giữ theo bản trích xuất.

\begin{verbatim}
#include <stdio.h>
#include <stdlib.h>
#include <string.h>

const int maxm=100+2;
const int maxn=100+2;

int n,m,dist[maxm],MaxT[maxm][maxn];
bool reachable[maxm][maxn],*can[maxm][maxm];
/*
m = so vu khi cua nuoc B.
n = so bom cua nuoc A.
dist[i] = neu bom A duoc phep dung lap lai, so bom it nhat de tieu diet
          cac vu khi B tu i den m.
MaxT[s][i] = bom i, neu bat dau no tu s, co the tieu diet den chi so lon nhat;
             neu bom i khong tieu diet duoc s thi MaxT[s][i]=s-1.
reachable[i][j] = bom j cua nuoc A co the tieu diet vu khi i cua nuoc B hay khong.
can[s][t][i] = bom i cua nuoc A co the tieu diet vu khi B tu s den t hay khong.
*/
int answer,bestv[maxn];
/*
answer = so bom A it nhat can dung.
bestv = ghi lai day bom A duoc dung trong nghiem toi uu.
*/
int a[maxm],b[maxn];
bool vis[maxn],*g[maxm];
/*
a[i], b[i] dung cho ghep cap; lan luot ghi dinh o phia con lai duoc ghep voi
dinh i ben trai/phai, neu chua ghep thi bang 0.
vis[i] dung de danh dau trong ghep cap va tim kiem theo chieu rong.
g[i][j] = co canh tu dinh i ben trai den dinh j ben phai hay khong.
*/
void init()
{
     // Doc du lieu va tinh reachable.
     int x1[maxm],y1[maxm],x2[maxn],y2[maxn],i,j,R;
     scanf("%d%d%d",&m,&n,&R);
     for (i=1;i<=m;i++)
        scanf("%d%d",&x1[i],&y1[i]);
     for (i=1;i<=n;i++)
        scanf("%d%d",&x2[i],&y2[i]);
     for (i=1;i<=m;i++)
        for (j=1;j<=n;j++)
           reachable[i][j]=
             ((x1[i]-x2[j])*(x1[i]-x2[j])+
              (y1[i]-y2[j])*(y1[i]-y2[j])<=R*R);
}

void preprocess()
{
    // Khoi tao, tinh can va MaxT.
    int s,t,i;
    for (i=1;i<=m;i++)
        g[i]=new bool[maxn];
    for (s=1;s<=m;s++)
        for (t=s;t<=m;t++)
            can[s][t]=new bool[maxn];
    for (s=1;s<=m;s++)
    {
        for (i=1;i<=n;i++)
            can[s][s][i]=reachable[s][i];
        for (t=s+1;t<=m;t++)
            for (i=1;i<=n;i++)
                 can[s][t][i]=can[s][t-1][i] && reachable[t][i];
        for (i=1;i<=n;i++)
        {
            MaxT[s][i]=s-1;
            for (t=s;t<=m;t++)
                 if (can[s][t][i])
                    MaxT[s][i]=t;
        }
    }
    // Tinh dist.
    dist[m+1]=0;
    for (s=m;s>=1;s--)
    {
        t=s-1;
        for (i=1;i<=n;i++)
            if (MaxT[s][i]>t)
                 t=MaxT[s][i];
        dist[s]=1+dist[t+1];
    }
}

bool find(int v)
{
    // Thuat toan Hungary tim duong tang.
    for (int i=1;i<=n;i++)
        if (g[v][i] && !vis[i])
        {
            vis[i]=true;
            if (b[i]==0 || find(b[i]))
            {
                a[v]=i;
                b[i]=v;
                return true;
            }
        }
    return false;
}

void search(int used,int s)
{
    // Trang thai: da dung used bom A; cac vu khi truoc s da bi pha huy.

    if (used+dist[s]>=answer) // Toi uu 1: cat tia toi uu.
       return;
    if (s==m+1)
    {
        // Tat ca vu khi B da bi pha huy: cap nhat nghiem toi uu roi quay lui.
        answer=used;
        memcpy(bestv,a,sizeof(a));
        return;
    }
    int t,maxL,tempa[maxm],tempb[maxn],op,cl,queue[maxn],k,i;
    // Tim kiem theo chieu rong de tinh do dai lon nhat maxL cua doan tiep theo.
    memset(vis,false,sizeof(vis));
    maxL=s-1;
    op=cl=0;
    for (i=1;i<=n;i++)
        if (b[i]==0)
        {
            vis[i]=true;
            queue[++op]=i;
        }
    while (cl<op)
    {
        k=queue[++cl];
        if (MaxT[s][k]>maxL)
            maxL=MaxT[s][k];
        for (i=1;i<=used;i++)
            if (g[i][k] && !vis[a[i]])
            {
                vis[a[i]]=true;
                queue[++op]=a[i];
            }
    }
    if (maxL==s-1)
        return;

    used++;
    memcpy(tempa,a,sizeof(a));
    memcpy(tempb,b,sizeof(b));
    memset(vis,false,sizeof(vis));
    g[used]=can[s][maxL];
    // Mo rong duong xen ke.
    find(used);
    // Liet ke do dai doan tiep theo tu lon den nho.
    for (t=maxL;t>=s;t--)
    {
        g[used]=can[s][t];
        search(used,t+1);
    }
    memcpy(a,tempa,sizeof(a));
    memcpy(b,tempb,sizeof(b));
}

void out()
{
    // In ket qua.
    printf("%d\n",answer);
    for (int i=1;i<=answer;i++)
        printf("%d ",bestv[i]);
    printf("\n");
}

int main()
{
    freopen("zplhz.in","r",stdin);
    freopen("zplhz.out","w",stdout);
    init();
    preprocess();
    answer=1000000000;
    memset(a,0,sizeof(a));
    memset(b,0,sizeof(b));
    search(0,1);
    out();
    return 0;
}
\end{verbatim}

\end{document}
